\section{Random-Function Coverage Baselines}
\label{sec:theory}

The concrete map $F_b$ in Eq.~\eqref{eq:iteration-map} is deterministic.
Classical TMTO analyses nevertheless model the underlying iteration function as
a random mapping on a set of size $N$.  This section states the random-function
baseline against which our exact DES-derived measurements are compared.

\subsection{Why independent occupancy is not a Hellman prediction}

If $q$ points were sampled independently and uniformly from an $N$-element
set, the expected number of distinct observed states would be
\[
    N\left(1-\left(1-\frac1N\right)^q\right)
    \approx N(1-e^{-q/N}).
\]
Putting $q=mt$ gives the familiar occupancy expression
\begin{equation}
    \rho_{\rm iid}
      \approx
      1-e^{-mt/N}.
    \label{eq:iid-occupancy}
\end{equation}
For $mt=N$, this is $1-e^{-1}\approx0.6321$.

Equation~\eqref{eq:iid-occupancy} is useful as a deliberately naive reference,
but Hellman matrix positions are not independent samples.  Within each row,
later entries are deterministic iterates of earlier entries; once two
trajectories meet under a fixed function, their suffixes coincide.  Classical
Hellman analysis and later random-mapping work explicitly account for this
dependence~\cite{Hellman1980,FlajoletOdlyzko1990,MaHong2009}.  We therefore do
not use Eq.~\eqref{eq:iid-occupancy} as the expected Hellman coverage.

\subsection{Ma--Hong expected coverage recurrence}

Ma and Hong~\cite{MaHong2009} analyze the first $t$ columns of an
$m\times(t+1)$ Hellman matrix.  Let
\[
   \overline H_k
   =
   \{H_{i,j}:1\leq i\leq m,\ 0\leq j\leq k\}
\]
and let $p_k$ denote, after normalization by $N$, the expected number of states
appearing for the first time in column $k$.  Introducing
\[
    s_k=\sum_{j=0}^{k-1}p_j,
\]
their large-$N$ approximation yields
\begin{equation}
    s_0=0,
    \qquad
    s_{k+1}
      =
      1-\exp\!\left(-\frac{m}{N}\right)\exp(-s_k).
    \label{eq:ma-hong-recurrence}
\end{equation}
The expected coverage rate relative to the $mt$ nominal matrix positions is
\begin{equation}
    \operatorname{ECR}(N,m,t)
      =
      \frac{N}{mt}s_t.
    \label{eq:ma-hong-ecr}
\end{equation}
Consequently, the predicted fraction of the \emph{whole state space} covered
by the first $t$ columns is simply
\begin{equation}
    \rho_{\rm MH}(N,m,t)
      =
      \frac{mt}{N}\operatorname{ECR}(N,m,t)
      =
      s_t.
    \label{eq:ma-hong-state-fraction}
\end{equation}

This distinction in normalization is particularly useful in our experiments.
The artifact reports $C_H/N$, whereas much of the TMTO literature reports
coverage relative to $mt$.

\subsection{Initial configuration}

For
\[
    N=2^{16}=65536,
    \qquad
    m=1024,
    \qquad
    t=64,
\]
Eq.~\eqref{eq:ma-hong-recurrence} gives
\[
    s_{64}\approx0.1665143030.
\]
Thus random-function theory predicts approximately
\[
    0.1665143030\,N
      \approx
      10912.68
\]
distinct represented states, or $16.6514\%$ of the state space.

The committed Java and Rust experiments report median exact Hellman coverage
near $16.8\%$.  The purpose of the larger sweep is to determine whether this
agreement persists across state dimensions and chain shapes, not merely at one
configuration.

For contrast, Eq.~\eqref{eq:iid-occupancy} gives
\[
    1-e^{-1}\approx63.2121\%
\]
because $mt=N$.  The difference between $63.2\%$ and $16.65\%$ is therefore
not itself a novel empirical anomaly; it is what established fixed-map
Hellman coverage theory already leads us to expect.

\subsection{The matrix-stopping parameter}

Hellman's standard parameterization is often discussed under the
\emph{matrix-stopping rule}
\begin{equation}
    mt^2\approx N,
    \label{eq:matrix-stopping}
\end{equation}
rather than $mt\approx N$~\cite{Hellman1980,MaHong2009}.  Under the exact
condition $mt^2=N$, Eq.~\eqref{eq:ma-hong-recurrence} becomes
\[
    s_{k+1}
      =
      1-\exp(-1/t^2)\exp(-s_k),
\]
and
\[
    \operatorname{ECR}(t)=t\,s_t.
\]

Our initial configuration instead satisfies
\[
    mt=N,
    \qquad
    mt^2=64N.
\]
It is therefore a useful high-overlap finite-state stress point, but it should
not be described as a standard matrix-stopping configuration.  The experiment
matrix includes both regimes:
\begin{enumerate}
  \item configurations near $mt^2=N$, to connect directly with classical
        parameterization;
  \item equal-$mt$ chain-shape sweeps, to isolate the effect of changing the
        number and length of rows at fixed nominal transition budget;
  \item the historical/repository stress configuration $m=1024,t=64$ at
        $N=2^{16}$.
\end{enumerate}

\subsection{Finite-$N$ and model limitations}

Equation~\eqref{eq:ma-hong-recurrence} is an approximation derived under a
random-function model.  It is not an exact identity for an arbitrary fixed
mapping.  Two sources of deviation are therefore relevant to the present
study.

First, the reduced state spaces are finite enough that the entire functional
graph can be enumerated.  Small-$N$ corrections, self-intersections, and
specific component structure need not equal their asymptotic expectations.
Ma and Hong themselves discuss a special discrepancy for the $m=1$ case and
give a separate exact expression~\cite{MaHong2009}.

Second, $F_b$ is not sampled uniformly from the set of all functions
$X_b\to X_b$: it is built from DES encryption, a fixed plaintext, a structured
key embedding, and a deterministic reduction.  Agreement with
Eq.~\eqref{eq:ma-hong-state-fraction} is therefore an empirical property to be
tested, not an assumption.

These observations motivate a direct random-map control.  For every state
dimension and TMTO parameter set used in the final experiments, we will compare
the DES-derived map with independently generated random mappings
\[
    F_{\rm rand}:X_b\longrightarrow X_b
\]
under the same start states and table parameters.

\subsection{Distinguished-point baseline}

For the predicate
\[
    D_d(x)=1
    \Longleftrightarrow
    x\bmod2^d=0,
\]
exactly $N/2^d$ states are distinguished.  If successive states behaved as
fresh uniform samples, the first distinguished hit would have geometric
parameter
\[
    p=2^{-d}.
\]
For an unbounded chain,
\[
    \Pr[\tau=k]=(1-p)^{k-1}p,
    \qquad
    \mathbb E[\tau]=2^d.
\]
With a hard transition bound $L$, the fresh-sample truncation probability is
\begin{equation}
    \Pr[\tau>L]
      =
      (1-2^{-d})^L.
    \label{eq:dp-truncation-baseline}
\end{equation}
This is again a baseline rather than an exact statement for iterations of a
fixed finite mapping.  Cycles that contain no distinguished state can make a
row fail to terminate even when $L$ is large, and merged trajectories create
dependence among rows.  The final experiments compare measured truncation and
coverage with Eq.~\eqref{eq:dp-truncation-baseline} and with exact
functional-graph structure.

\subsection{What constitutes a meaningful deviation}

For each configuration we compute the theory residual
\[
    \Delta_{\rm cov}
      =
      \rho_{\rm measured}-\rho_{\rm MH}.
\]
A single nonzero residual is not evidence of cryptographic structure:
finite-size variability and the approximation itself must be accounted for.
The final paper therefore reports distributions over repeated random start
sets and, for the random-map controls, over independently sampled mappings.
Only a deviation that is stable relative to these controls will be interpreted
as evidence that the DES-derived map/reduction has coverage behavior not
captured by the random-function model.
